\chapter{Critérios de Seleção de Artigos e Sensores}
\label{chap:appendix-selection-criteria}

\begin{introduction}
Este apêndice detalha o processo e critérios utilizados na seleção de artigos, bem como na seleção de fontes de dados a mencionar na Secção~\ref{sec:soa_data_sources}
\end{introduction}

\section{Metodologia de Pesquisa de Artigos Científicos}

A pesquisa bibliográfica de artigos foi conduzida principalmente através da plataforma \textit{Google Scholar}, utilizando um conjunto de palavras-chave diferente conforme necessário. Esta pesquisa inicial foi posteriormente complementada através da identificação de referências relevantes presentes em artigos já selecionados.

De modo a manter a relevância temporal dos artigos, e tendo em conta a constante evolução da tecnologia, foram priorizados artigos publicados após 2022, existindo, contudo, exceções para artigos cuja qualidade, importância ou nível de detalhe o justifiquem. 
Como critério de qualidade mínima, avaliei também a reputação e processo de seleção da publicadora ou plataforma em que o artigo se encontra, já que critérios de seleção exigentes tendem a refletir-se na qualidade dos artigos aceites.

Finalmente, as citações utilizadas não correspondem a uma lista exaustiva de todos os resultados válidos nem necessariamente aos ``melhores'' artigos da área, mas apenas a exemplos representativos que adequadamente sustentam a afirmação.

\section{Critérios de Seleção de Fontes de Dados}

Na Secção~\ref{sec:soa_data_sources} foram listadas as principais fontes de dados encontradas nos artigos para monitorização de utilizadores.
Devido à especificidade do contexto, foram incluídos sensores de monitorização mesmo quando estes não foram explicitamente utilizados enquanto o indivíduo interagia com o computador, mas apenas se for considerado que estes não impeçam fisicamente os movimentos necessários para a interação com o computador.
Por esta razão, não foi mencionado o uso do \textit{smartphone} e dos sensores internos deste nesta secção, já que a captação dos dados exige geralmente que o utilizador segure nele e o manipule, impedindo-o de realizar tarefas como digitação.